\newpage
\section{Analysis Implementations}

a brief overview here about what will be covered in this section.

\subsection{Mean Squared Displacement}

This is a section that goes over the mathematical background that is followed
or was followed while implementing the code. Mean squared displacement is a 
quantity that allows us to see how far a particle's position deviates from a 
reference point over time. The MSD at a time $t$ is defined as the following
ensemble average:

\begin{equation*}
    MSD = \langle| \mathbf{x}(t) - \mathbf{x}(0)|\rangle = \frac{1}{N} \sum_{i=1}^N 
    |\mathbf{x}^{(i)}(t) - x^{(i)}(0)|^2.
\end{equation*}

Maybe also mention something in this section about MSD for different lags. Also
maybe discuss what we should see/what MSD should look like and what the 
characteristics of a specific plot mean (the inflection point, etc.).



\subsection{Radial Distribution Function}

This section should go over the mathematical background of the radial
distribution function calculation. Specifically, what expression/whatever is 
actually used in the code to calculate the RDF. Maybe in these sections or in 
the implementation sections, we could put in addition to the actual code, just
some pseudo-code for thinking about how the code should be implemented in the
first place. Then, the rest is just a C++-specific implementation.

Again, also describe here certain features in the plots (should go to 1 for
large distances, what each peak represents, etc.)

The radial distribution or pair correlation function, describes how density
varies as a function of distance from some reference particle. It is given by:

EQUATION/CALCULATION HERE


\subsection{First Passage Times}

Don't forget to add citations for Jerome's paper in this section.

The mean first passage time is given by:
\begin{equation*}
    \mathcal{F}(r) = \int_{0}^{\infty} dt t f(r, t) = - \frac{\partial \hat{f}(r,s)}{\partial s}
\end{equation*}
We also have a quantity, the logarithmic derivative of this quantity, which we're 
using to make the quantity or see the feature ...

\begin{equation*}
    D(r) = \frac{d \ln{\mathcal{F}(r)}}{d \ln{r}}.
\end{equation*}


\subsection{Correlation Functions}

\subsubsection{Velocity Autocorrelation Functions}

\subsubsection{Current Correlation Functions}

\subsubsection{Structure Factors}

\subsubsection{Intermediate Scattering}

\subsubsection{Custom Correlations}
